\documentclass[10pt,a4paper]{amsart}
\usepackage[margin=19mm]{geometry}
\usepackage{amsmath,amssymb,mathtools}
\usepackage[scr=rsfs,cal=euler]{mathalfa}
\usepackage{xfrac}
\usepackage[unicode,hidelinks,bookmarks=false]{hyperref}
\newcommand{\ie}{i.e.\ }
\newcommand{\cf}{cf.\ }
\newcommand{\resp}{respectively\ }
\newcommand{\Subsection}{\subsection}
\providecommand{\nobreakdash}{\nobreak\hbox{-}}
\setlength{\emergencystretch}{2em}
\multlinegap0pt
\usepackage{enumerate}
\usepackage{amssymb}
\usepackage{xcolor}
\usepackage{dsfont}
\usepackage[normalem]{ulem}
\newtheorem{lemma}{Lemma}
\usepackage{cleveref}
\crefname{lemma}{Lemma}{Lemmas}
\newcommand{\Gauss}[2]{{\begin{bmatrix} #1 \\ #2 \end{bmatrix}_q}}
\newcommand{\PS}{\mathrm{PS}}
\renewcommand{\geq}{\geqslant}
\renewcommand{\leq}{\leqslant}
\title{The largest sets of non-opposite chambers in spherical buildings of type $B$}
\author{Jan De Beule, Philipp Heering, Sam Mattheus, Klaus Metsch}
\date{Source: arXiv:2505.14322v2 (CC BY 4.0)}
\begin{document}
\maketitle

\noindent\emph{Source and licence.} The largest sets of non-opposite chambers in spherical buildings of type $B$. Version arXiv:2505.14322v2 (CC BY 4.0), distributed by arXiv under the Creative Commons Attribution 4.0 International licence, http://creativecommons.org/licenses/by/4.0/, as recorded in the arXiv licence metadata for this exact version and shown on the abstract page https://arxiv.org/abs/2505.14322v2. Full licence notice: \texttt{licenses/arxiv-2505.14322-CC-BY-4.0.txt}.\par\smallskip

\noindent\textit{Editorial scope.} Counted unit: the article's lemma perp (source0.tex lines 403--405). The article's complete original proof of it is delivered with it (source0.tex lines 406--440, one \texttt{proof} environment of 35 lines). No mathematics is written, repaired or re-derived: the statement and the proof are the article's own lines, in the article's own notation, and no retained line is substituted or re-flowed.  Retained uncounted as the source-local context the proof needs: lines 319--327.  Nothing was omitted from any retained span, so the package carries no omission record.  No retained line contains a citation, so the wrapper carries no bibliography and no bibliography entry.  No retained line contains a figure environment, an image-inclusion command or drawing code, so no image is delivered and the package declares no asset dependency.  Editorial formatting only, in addition: an AMS \texttt{amsart} preamble that restates the article's own declarations and macros used by the retained lines, namely the environments \texttt{lemma} on separate counters, together with the standard packages those lines need.  The retained environments print in this document as Lemma 1 in order of appearance, whereas the article prints the same items with the numbering of its own full document.  The complete native source of the article as distributed in the arXiv e-print for this exact version is bundled unchanged as \texttt{sources/178431/source0.tex}.
\begin{lemma} \label{L:degree_stuff}
For $1\leq s\leq n$, we have
    \begin{enumerate} [(a)]
        \item $\deg(z_s)=s(s-1)/2$,
        \item $\deg(\Phi_m^s)=(s-m)(2n-s-m+e-\frac{s-m+1}{2})$,
        \item $\deg(z_sz_{n-s}\Phi_s^n)=s(s-1)/2+(n-s)(n-s+e-1)$,
        \item $\deg(\Phi_1^nz_n)=(n-1)(n+e-1)$.
    \end{enumerate}
\end{lemma}

\begin{lemma} \label{L: points in X perp}
    Let $x$ be a real number and let $X$ be a non-empty set of $\Omega(q^{x-1})$ points in $\PS(n,e,q)$, where $n\geq 2$ and $e\geq 1$. Then $\langle X \rangle^\perp$ contains $O(q^{N-x})$ (singular) points, where $N= \deg(\Phi_0^1(n,e,q))$. 
\end{lemma}

\begin{proof}
According to \Cref{L:degree_stuff}, we have $N=2n+e-2$.
As $X$ contains  $\Omega(q^{x-1})$ points, we have $\dim(\langle X\rangle)\geq x$.
   We consider the cases independently.\\

    \underline{Case 1: symplectic polar spaces.} We have a symplectic polar space, so $e=1$ and the ambient vector space of $\PS(n,e,q)$ has dimension $2n$. As the alternating form is non-degenerate, we have $\dim(\langle X \rangle^\perp)\leq 2n-x=N+1-x$, which yields that $\langle X \rangle^\perp$ contains $O(q^{N-x})$ singular points.  \\
    
    \underline{Case 2: quadratic or Hermitian spaces.}
    These are the cases where $\PS(n,e,q)$ is parabolic with odd characteristic ($e=1$), elliptic ($e=2$), or Hermitian ($e=3/2$).\\  
    
    Assume that the underlying quadratic or Hermitian form does not vanish identically on $\langle X \rangle^\perp$.
    In the quadratic case, the ambient vector space has dimension $2n+e$. Since the associated polar form is non-degenerate, we have $\dim(\langle X \rangle^\perp)\leq 2n+e-x$. The singular points of $\langle X \rangle^\perp$ form the zero set of a non-zero quadratic form. Hence, $|\langle X \rangle^\perp \cap\PS(n,e,q)|= O(q^{2n+e-x-2})
    =O(q^{N-x})$.\\
    In the Hermitian case, the ambient vector space has dimension $2n+1$. Since the Hermitian form is non-degenerate, we have $\dim(\langle X \rangle^\perp)\leq 2n+1-x$. The singular points of $\langle X \rangle^\perp$ form the zero set of a non-zero Hermitian form. Therefore, $|\langle X \rangle^\perp\cap\PS(n,e,q)|= O(q^{2n-x-1/2})=O(q^{N-x})$.\\
    
Now, suppose that the defining quadratic or Hermitian form vanishes on $\langle X\rangle^\perp$. Then $\langle X\rangle^\perp$ is totally singular, and hence $|\langle X\rangle^\perp\cap\PS(n,e,q)|=O(q^{\dim(\langle X\rangle^\perp)-1})$.
As the points of $\PS(n,e,q)$ contained in $\langle X\rangle$ form a cone whose radical is $\langle X\rangle^\perp$ and whose base is a polar space of rank $n-\dim(\langle X\rangle^\perp)$, we have 
\begin{align} \label{Align: 2.10}
    |\langle X\rangle\cap\PS(n,e,q)|&=\Gauss{\dim(\langle X\rangle^\perp)}{1}+
    q^{\dim(\langle X\rangle^\perp)}
    \Phi_0^1(n-\dim(\langle X\rangle^\perp),e,q). %\\
   % &=
   % O(q^{N-\dim(\langle X\rangle^\perp)}).
\end{align}
%The term $\Gauss{\dim(\langle X\rangle^\perp)}{1}$ is $O(q^{\dim(\langle X\rangle^\perp)-1})$.
First, let us consider the case $\dim(\langle X\rangle^\perp)<n$.
We can use \Cref{L:degree_stuff} (b) and obtain that $\Phi_0^1(n-\dim(\langle X\rangle^\perp),e,q)=O(q^{2n-2\dim(\langle X\rangle^\perp)+e-2})=O(q^{N-2\dim(\langle X\rangle^\perp)})$. So, the first summand of (\ref{Align: 2.10}) is  $O(q^{\dim(\langle X\rangle^\perp)-1})$ and the second summand of (\ref{Align: 2.10}) is $O(q^{N-\dim(\langle X\rangle^\perp)})$. As $\dim(\langle X\rangle^\perp)-1<N-\dim(\langle X\rangle^\perp)$, we get $|\langle X\rangle\cap\PS(n,e,q)|=O(q^{N-\dim(\langle X\rangle^\perp)})$.\\
Now, let us consider the case $\dim(\langle X\rangle^\perp)=n$. Then $\Phi_0^1(n-\dim(\langle X\rangle^\perp),e,q)=0$.  Hence, we get $|\langle X\rangle\cap\PS(n,e,q)|=O(q^{\dim(\langle X\rangle^\perp)-1})$. As $n-1\leq N-n$ for $e\geq 1$, we obtain $|\langle X\rangle\cap\PS(n,e,q)|=O(q^{N-\dim(\langle X\rangle^\perp)})$ in both cases.

We have $X\subseteq \langle X \rangle\cap \PS(n,e,q)$ and as $|X|=\Omega(q^{x-1})$, this gives $x-1\leq N-\dim(\langle X\rangle^\perp)$, which implies the statement.\\

\underline{Case 3: parabolic quadrics in even characteristic.}
Let $R$ be the one-dimensional radical of the associated polar bilinear form. Projection from $R$ induces an isomorphism to a symplectic polar space, which preserves cardinalities and common perpendiculars.
Therefore, the statement follows from Case 1.
\end{proof}

\end{document}
